
\documentclass[11pt,a4paper]{article}

\usepackage[utf8]{inputenc}
\usepackage[T1]{fontenc}
\usepackage[spanish,es-nodecimaldot]{babel}
\usepackage{amsmath,amssymb,amsthm,mathtools}
\usepackage{bm}
\usepackage{geometry}
\usepackage{booktabs}
\usepackage{array}
\usepackage{enumitem}
\usepackage{hyperref}
\usepackage{xcolor}
\usepackage{microtype}
\usepackage{graphicx}
\usepackage{longtable}
\usepackage{listings}
\usepackage{fancyhdr}
\usepackage{titlesec}

\geometry{margin=2.4cm}
\hypersetup{
    colorlinks=true,
    linkcolor=blue!60!black,
    citecolor=blue!60!black,
    urlcolor=blue!60!black
}

\pagestyle{fancy}
\fancyhf{}
\lhead{Nota de trabajo --- AlphaZero}
\rhead{Pedro E. Colla}
\cfoot{\thepage}

\title{\textbf{AlphaZero: fundamentos, formulación matemática, búsqueda, aprendizaje y construcción de estrategia}\\
\large Nota de trabajo técnica}
\author{Dr. Pedro E. Colla}
\date{\today}

\newtheorem{definition}{Definición}[section]
\newtheorem{proposition}{Proposición}[section]
\newtheorem{remark}{Observación}[section]

\newcommand{\E}{\mathbb{E}}
\newcommand{\R}{\mathbb{R}}
\newcommand{\Pp}{\mathbb{P}}
\newcommand{\argmax}{\operatorname*{arg\,max}}
\newcommand{\argmin}{\operatorname*{arg\,min}}
\newcommand{\softmax}{\operatorname{softmax}}
\newcommand{\KL}{D_{\mathrm{KL}}}

\begin{document}

\maketitle
\tableofcontents
\newpage

\begin{abstract}
AlphaZero es un algoritmo de aprendizaje por refuerzo profundo que aprende a jugar juegos de información perfecta y suma cero a partir de autojuego, sin utilizar partidas humanas como fuente necesaria de supervisión. Su arquitectura combina una red neuronal que estima simultáneamente una política y un valor, con una búsqueda Monte Carlo Tree Search (MCTS) guiada mediante una regla de exploración tipo PUCT. La búsqueda produce una política mejorada respecto de la política cruda de la red; el autojuego genera datos; y esos datos se utilizan para entrenar una nueva red. El proceso se repite iterativamente.

Esta nota desarrolla de forma autocontenida los conceptos matemáticos y algorítmicos necesarios para comprender el funcionamiento de AlphaZero, incluyendo la formulación como proceso de decisión de Markov, la representación del estado, la red de política-valor, la búsqueda PUCT, el autojuego, la función de pérdida, la actualización de parámetros, la relación con iteración de políticas generalizada y las limitaciones teóricas del método. El objetivo es distinguir con precisión entre lo que AlphaZero calcula, lo que aproxima y lo que puede inferirse de sus resultados.
\end{abstract}

\section{Motivación general}

AlphaZero pertenece a la familia de algoritmos que combinan aprendizaje por refuerzo, redes neuronales profundas y búsqueda. Su idea central puede resumirse como un ciclo de mejora:

\begin{enumerate}[label=\arabic*.]
    \item una red neuronal propone qué jugadas parecen prometedoras y estima el valor de una posición;
    \item una búsqueda de árbol utiliza esas estimaciones para analizar selectivamente variantes;
    \item la distribución de visitas del árbol se interpreta como una política mejorada;
    \item se juega utilizando esa política;
    \item el resultado final de la partida se propaga como señal de aprendizaje;
    \item la red se entrena para aproximar las políticas mejoradas y los resultados observados;
    \item el procedimiento se repite.
\end{enumerate}

Esto crea una interacción entre dos mecanismos distintos:

\begin{itemize}
    \item \textbf{generalización}: la red neuronal intenta capturar regularidades globales;
    \item \textbf{razonamiento local}: MCTS examina con mayor detalle la posición concreta.
\end{itemize}

La red sola es rápida pero imperfecta. La búsqueda sola sería muy costosa en espacios enormes. AlphaZero combina ambas.

\section{Marco matemático: juegos como procesos de decisión}

\subsection{Estados, acciones y transición}

Sea un juego determinista de dos jugadores, de información perfecta y suma cero. Denotemos:

\[
s \in \mathcal{S}
\]

al estado del juego, y

\[
a \in \mathcal{A}(s)
\]

a una acción legal en el estado \(s\).

La transición determinista se escribe:

\[
s' = T(s,a).
\]

En ajedrez, por ejemplo, \(s\) contiene como mínimo:

\begin{itemize}
    \item disposición de piezas;
    \item jugador al turno;
    \item derechos de enroque;
    \item posibilidad de captura al paso;
    \item información relevante para reglas de tablas.
\end{itemize}

\subsection{Recompensa terminal}

En la formulación original de AlphaZero, la señal principal de recompensa aparece al final de la partida. Si \(z\) representa el resultado desde la perspectiva del jugador correspondiente, típicamente:

\[
z =
\begin{cases}
+1, & \text{victoria},\\
0, & \text{tablas},\\
-1, & \text{derrota}.
\end{cases}
\]

El juego es de suma cero, de modo que si para un jugador el resultado es \(+1\), para el otro es \(-1\).

\subsection{Política}

Una política es una distribución de probabilidad sobre acciones legales:

\[
\pi(a\mid s).
\]

Satisface:

\[
\pi(a\mid s)\ge 0,
\qquad
\sum_{a\in \mathcal{A}(s)}\pi(a\mid s)=1.
\]

La política puede verse como una regla probabilística de decisión.

\subsection{Función de valor}

El valor de un estado bajo una política \(\pi\) es:

\[
V^\pi(s)
=
\E_\pi[z\mid s].
\]

En un juego de suma cero, un valor positivo indica ventaja esperada para el jugador cuya perspectiva se está utilizando.

Una política óptima, idealmente, satisface:

\[
V^*(s)
=
\max_\pi V^\pi(s),
\]

pero AlphaZero no calcula exactamente esta cantidad en juegos grandes. La aproxima mediante red neuronal y búsqueda.

\section{El problema de ``la mejor estrategia''}

Es importante distinguir tres objetos diferentes:

\begin{enumerate}
    \item la política cruda de la red neuronal;
    \item la política refinada por MCTS;
    \item la política verdaderamente óptima del juego.
\end{enumerate}

La red produce:

\[
p_\theta(a\mid s),
\]

donde \(\theta\) son los parámetros de la red.

La búsqueda produce otra distribución:

\[
\pi_{\mathrm{MCTS}}(a\mid s).
\]

La política efectiva de AlphaZero durante el autojuego se deriva principalmente de las visitas del árbol, no simplemente del máximo de \(p_\theta\).

La política óptima real:

\[
\pi^*(a\mid s)
\]

es, en juegos enormes, desconocida.

Por tanto:

\[
p_\theta
\;\neq\;
\pi_{\mathrm{MCTS}}
\;\neq\;
\pi^*
\]

en general.

AlphaZero intenta que, a medida que el entrenamiento progresa,

\[
p_\theta \approx \pi_{\mathrm{MCTS}},
\]

y que el proceso conjunto acerque ambas políticas a decisiones de alta calidad.

\section{La red neuronal de política y valor}

\subsection{Función conjunta}

La red neuronal implementa una función:

\[
f_\theta(s) = (\bm{p}, v),
\]

donde:

\[
\bm{p} = \left(p(a_1\mid s),\ldots,p(a_n\mid s)\right)
\]

es una distribución sobre las acciones posibles y

\[
v \in [-1,1]
\]

es la estimación del resultado esperado desde el estado \(s\).

Esto implica que una sola red responde dos preguntas:

\begin{enumerate}
    \item \textbf{Política:} ¿qué jugadas parecen prometedoras?
    \item \textbf{Valor:} ¿qué tan favorable parece la posición?
\end{enumerate}

\subsection{Cabeza de política}

Internamente la red produce logits \(l_a\). La política se obtiene mediante softmax:

\[
p(a\mid s)
=
\frac{e^{l_a}}
{\sum_b e^{l_b}}.
\]

La política no es todavía la decisión final. Funciona como \emph{prior} para orientar la búsqueda.

\subsection{Cabeza de valor}

La cabeza de valor produce un escalar:

\[
v_\theta(s)
\approx
\E[z\mid s].
\]

Normalmente se restringe a:

\[
v_\theta(s)\in[-1,1],
\]

por ejemplo mediante una activación \(\tanh\).

\subsection{Representación del tablero}

En ajedrez, el estado se codifica como tensores espaciales. Conceptualmente:

\[
X(s)\in\R^{H\times W\times C}.
\]

Para un tablero de ajedrez:

\[
H=W=8,
\]

y \(C\) contiene planos que representan piezas, turno, historial y otra información.

La ventaja de esta representación es preservar la estructura espacial del tablero.

\section{Arquitectura residual}

AlphaZero utiliza una red profunda con bloques residuales. Si \(x_l\) es la entrada a un bloque, la forma conceptual es:

\[
x_{l+1}
=
\sigma\left(
x_l + F(x_l;W_l)
\right).
\]

La conexión residual:

\[
x_l \rightarrow x_l + F(x_l)
\]

facilita el entrenamiento de redes profundas.

El cuerpo compartido de la red extrae características; luego se divide en:

\[
\text{trunk}
\longrightarrow
\begin{cases}
\text{policy head},\\
\text{value head}.
\end{cases}
\]

\section{Monte Carlo Tree Search}

\subsection{Idea básica}

MCTS construye un árbol de posiciones alcanzables. Cada arista \((s,a)\) almacena estadísticas tales como:

\[
N(s,a): \text{número de visitas},
\]

\[
W(s,a): \text{suma de valores acumulados},
\]

\[
Q(s,a): \text{valor medio},
\]

\[
P(s,a): \text{prior de la red}.
\]

Con:

\[
Q(s,a)
=
\frac{W(s,a)}{N(s,a)}
\]

si \(N(s,a)>0\).

\subsection{Las cuatro fases conceptuales}

Una simulación MCTS contiene cuatro etapas:

\begin{enumerate}
    \item selección;
    \item expansión;
    \item evaluación;
    \item retropropagación.
\end{enumerate}

Estas etapas se repiten muchas veces desde la misma posición raíz.

\section{Selección mediante PUCT}

\subsection{Regla básica}

AlphaZero utiliza una forma de Upper Confidence Bound guiada por una política previa. Una expresión típica es:

\[
a_t
=
\argmax_a
\left[
Q(s,a)
+
U(s,a)
\right],
\]

donde

\[
U(s,a)
=
c_{\mathrm{puct}}\,
P(s,a)
\frac{\sqrt{\sum_b N(s,b)}}
{1+N(s,a)}.
\]

Por tanto:

\[
a_t
=
\argmax_a
\left[
Q(s,a)
+
c_{\mathrm{puct}}\,
P(s,a)
\frac{\sqrt{\sum_b N(s,b)}}
{1+N(s,a)}
\right].
\]

\subsection{Interpretación}

El término:

\[
Q(s,a)
\]

favorece jugadas que ya han producido buenos resultados.

El término:

\[
U(s,a)
\]

favorece jugadas:

\begin{itemize}
    \item poco visitadas;
    \item consideradas prometedoras por la red;
    \item relevantes en relación con el número total de visitas del nodo.
\end{itemize}

Esta tensión se denomina:

\[
\text{exploración}
\quad\text{vs}\quad
\text{explotación}.
\]

\subsection{Efecto de \(c_{\mathrm{puct}}\)}

Si \(c_{\mathrm{puct}}\) es grande, la búsqueda explora más.

Si \(c_{\mathrm{puct}}\) es pequeño, la búsqueda explota antes las variantes que ya parecen mejores.

No existe un valor universalmente óptimo; su elección depende de la escala de la red, el juego y el presupuesto de simulaciones.

\section{Expansión de un nodo}

Cuando la búsqueda alcanza por primera vez un estado \(s_L\), la red se evalúa:

\[
(\bm{p},v)
=
f_\theta(s_L).
\]

Para cada acción legal \(a\):

\[
P(s_L,a)
=
p_\theta(a\mid s_L).
\]

El valor:

\[
v_\theta(s_L)
\]

se utiliza como estimación del resultado de la posición sin tener que jugar una simulación aleatoria hasta el final.

Ésta es una diferencia importante respecto del MCTS clásico basado en rollouts aleatorios.

\section{Retropropagación}

Supongamos que una simulación alcanza un estado hoja con evaluación:

\[
v.
\]

Ese valor se propaga por las aristas visitadas.

Para una arista:

\[
N(s,a)
\leftarrow
N(s,a)+1,
\]

\[
W(s,a)
\leftarrow
W(s,a)+v,
\]

\[
Q(s,a)
\leftarrow
\frac{W(s,a)}{N(s,a)}.
\]

Debido a la alternancia de jugadores, debe corregirse la perspectiva del valor. Si un nodo se evalúa desde el jugador que mueve allí, al retroceder un nivel:

\[
v \leftarrow -v.
\]

Esta alternancia de signo es fundamental en juegos de suma cero.

\section{De las visitas del árbol a una política mejorada}

Después de \(M\) simulaciones, cada acción raíz posee un número de visitas:

\[
N(s,a).
\]

AlphaZero construye una política de búsqueda:

\[
\pi(a\mid s)
=
\frac{N(s,a)^{1/\tau}}
{\sum_b N(s,b)^{1/\tau}},
\]

donde \(\tau\) es un parámetro de temperatura.

\subsection{Temperatura alta}

Para \(\tau\) relativamente grande, la distribución es más plana:

\[
\pi(a\mid s)
\approx
\text{exploratoria}.
\]

\subsection{Temperatura baja}

Cuando:

\[
\tau\rightarrow 0,
\]

la política se aproxima a seleccionar:

\[
a^*
=
\argmax_a N(s,a).
\]

En entrenamiento suele utilizarse mayor exploración en las primeras jugadas y decisiones más deterministas después.

\section{Ruido de Dirichlet}

Para evitar que el autojuego se vuelva prematuramente determinista, AlphaZero perturba los priors en la raíz:

\[
P'(s,a)
=
(1-\varepsilon)P(s,a)
+
\varepsilon \eta_a,
\]

donde

\[
\bm{\eta}
\sim
\operatorname{Dir}(\alpha).
\]

Esto induce exploración estructurada.

El propósito no es mejorar directamente una jugada individual, sino incrementar la diversidad del conjunto de experiencias.

\section{El ciclo de autojuego}

Una partida de autojuego produce una secuencia:

\[
s_0,a_0,s_1,a_1,\ldots,s_T.
\]

Para cada estado \(s_t\):

\begin{enumerate}
    \item se ejecuta MCTS;
    \item se obtiene \(\pi_t\);
    \item se elige una acción según \(\pi_t\);
    \item se avanza al siguiente estado.
\end{enumerate}

Al finalizar la partida se obtiene:

\[
z\in\{-1,0,+1\}.
\]

Cada posición se transforma en un ejemplo:

\[
(s_t,\pi_t,z_t).
\]

El conjunto de datos de entrenamiento es entonces:

\[
\mathcal{D}
=
\{(s_i,\pi_i,z_i)\}_{i=1}^{N}.
\]

\section{El objetivo de aprendizaje}

La red debe aprender simultáneamente:

\begin{enumerate}
    \item a reproducir la política mejorada de MCTS;
    \item a predecir el resultado final.
\end{enumerate}

Una función de pérdida típica es:

\[
\mathcal{L}(\theta)
=
(z-v_\theta(s))^2
-
\bm{\pi}^{\top}
\log \bm{p}_\theta(s)
+
\lambda
\|\theta\|_2^2.
\]

\subsection{Pérdida de valor}

\[
\mathcal{L}_V
=
(z-v)^2.
\]

Hace que:

\[
v_\theta(s)\rightarrow z.
\]

\subsection{Pérdida de política}

\[
\mathcal{L}_P
=
-
\sum_a
\pi(a\mid s)
\log p_\theta(a\mid s).
\]

Es la entropía cruzada entre la política producida por MCTS y la política de la red.

Equivale a minimizar, salvo una constante:

\[
\KL(\pi\parallel p_\theta).
\]

\subsection{Regularización}

\[
\mathcal{L}_R
=
\lambda\|\theta\|_2^2.
\]

Su función es reducir el sobreajuste y limitar el crecimiento excesivo de parámetros.

\section{Actualización por descenso de gradiente}

Los parámetros se ajustan mediante:

\[
\theta_{k+1}
=
\theta_k
-
\eta
\nabla_\theta
\mathcal{L}(\theta_k),
\]

donde:

\[
\eta
\]

es la tasa de aprendizaje.

En la práctica se utilizan optimizadores por mini-batches y variantes de descenso estocástico.

\section{Por qué la búsqueda puede ser mejor que la red}

La red devuelve una evaluación local:

\[
f_\theta(s)=(p,v).
\]

Pero MCTS construye información adicional mediante múltiples expansiones.

En términos conceptuales:

\[
p_\theta
\overset{\mathrm{MCTS}}{\longrightarrow}
\pi.
\]

La búsqueda puede descubrir que una jugada inicialmente poco probable conduce a posiciones fuertes, o que una jugada aparentemente atractiva tiene respuestas adversas.

Por eso:

\[
\pi
\]

se considera un objetivo mejorado respecto de:

\[
p_\theta.
\]

\section{AlphaZero como iteración de políticas generalizada}

El procedimiento puede interpretarse como una forma aproximada de \emph{Generalized Policy Iteration}.

Sea:

\[
\pi_k
\]

la política representada por la red en una iteración.

MCTS produce:

\[
\pi_k'
=
\operatorname{Search}(\pi_k,V_k).
\]

Luego el entrenamiento ajusta:

\[
\pi_{k+1}
\approx
\pi_k'.
\]

Simultáneamente:

\[
V_{k+1}
\approx
V^{\pi_k'}.
\]

El ciclo conceptual es:

\[
(\pi_k,V_k)
\rightarrow
\mathrm{MCTS}
\rightarrow
(\pi_k',z)
\rightarrow
\mathrm{entrenamiento}
\rightarrow
(\pi_{k+1},V_{k+1}).
\]

\section{Relación con Bellman}

La ecuación óptima de Bellman para un problema de decisión es:

\[
V^*(s)
=
\max_a
\left[
r(s,a)
+
\gamma
V^*(T(s,a))
\right].
\]

En un juego determinista de suma cero, la alternancia de jugadores introduce un cambio de signo. Una forma simplificada puede escribirse:

\[
V^*(s)
=
\max_a
\left[
- V^*(T(s,a))
\right]
\]

para estados no terminales sin recompensa inmediata.

AlphaZero no aplica una tabla Bellman explícita sobre cada estado. En cambio:

\begin{itemize}
    \item la red aproxima el valor;
    \item MCTS realiza una forma de respaldo de valores;
    \item el resultado terminal proporciona la señal de entrenamiento.
\end{itemize}

Por ello existe una conexión conceptual profunda con Bellman, pero el mecanismo computacional es distinto de una tabla Q clásica.

\section{Valor de una acción}

Aunque la red produce principalmente \(V(s)\), MCTS mantiene:

\[
Q(s,a).
\]

Una interpretación es:

\[
Q(s,a)
\approx
- V(T(s,a))
\]

desde la perspectiva del jugador que mueve en \(s\), aunque en la práctica \(Q\) se obtiene mediante promedios de simulaciones y evaluaciones descendientes.

\section{Cómo se define una jugada ``mejor''}

En una posición concreta, AlphaZero no utiliza un único criterio.

La decisión surge de una combinación de:

\[
P(s,a),
\quad
Q(s,a),
\quad
N(s,a),
\quad
U(s,a).
\]

Durante la búsqueda:

\[
a_{\mathrm{select}}
=
\argmax_a
\left[
Q(s,a)+U(s,a)
\right].
\]

Al finalizar la búsqueda, la elección suele basarse en:

\[
N(s,a).
\]

Por tanto, la jugada final no es necesariamente:

\[
\argmax_a P(s,a),
\]

ni tampoco:

\[
\argmax_a Q(s,a).
\]

Es el resultado acumulado del proceso de búsqueda.

\section{Política de entrenamiento versus política de juego}

Durante entrenamiento:

\[
a_t\sim \pi_t.
\]

Esto genera exploración.

Durante evaluación competitiva puede utilizarse:

\[
a_t
=
\argmax_a N(s_t,a).
\]

Esto elimina gran parte de la aleatoriedad.

Esta distinción es esencial:

\[
\text{explorar para aprender}
\neq
\text{explotar para competir}.
\]

\section{Aprendizaje desde cero}

Una propiedad distintiva del enfoque AlphaZero es que no requiere, en principio:

\begin{itemize}
    \item aperturas humanas;
    \item bases de finales humanas;
    \item heurísticas posicionales diseñadas manualmente;
    \item partidas de grandes maestros.
\end{itemize}

Su conocimiento emerge de:

\[
\text{reglas}
+
\text{autojuego}
+
\text{búsqueda}
+
\text{optimización}.
\]

Esto no significa que ``no haya conocimiento incorporado''. La arquitectura, representación, reglas, parámetros de búsqueda y procedimiento de entrenamiento constituyen fuertes sesgos inductivos.

\section{Bootstrap y problema inicial}

Al comienzo:

\[
p_\theta(a\mid s)
\approx
\text{ruido estructurado},
\]

\[
v_\theta(s)
\approx
0.
\]

Sin embargo, el sistema puede comenzar a aprender porque:

\begin{enumerate}
    \item juega partidas;
    \item obtiene resultados terminales reales;
    \item utiliza esos resultados como objetivos de valor;
    \item utiliza MCTS para producir objetivos de política;
    \item las mejoras pequeñas se retroalimentan.
\end{enumerate}

Este proceso es un \emph{bootstrap}.

\section{Replay buffer}

Los ejemplos generados se almacenan normalmente en un buffer:

\[
\mathcal{B}
=
\{(s,\pi,z)\}.
\]

El entrenamiento extrae mini-batches:

\[
B
\subset
\mathcal{B}.
\]

Esto desacopla parcialmente:

\[
\text{generación}
\quad\text{y}\quad
\text{aprendizaje}.
\]

Un riesgo importante es entrenar demasiado sobre experiencias obsoletas. Si la política cambia rápidamente, un buffer excesivamente viejo puede contener objetivos \(\pi\) producidos por redes mucho más débiles.

\section{Distribución cambiante de datos}

AlphaZero aprende de una distribución que él mismo genera:

\[
s
\sim
d_{\pi_\theta}.
\]

Cuando cambia la política:

\[
\pi_\theta
\rightarrow
\pi_{\theta'},
\]

también cambia la distribución:

\[
d_{\pi_\theta}
\rightarrow
d_{\pi_{\theta'}}.
\]

El problema de aprendizaje es entonces no estacionario.

Esto es una de las razones por las cuales la estabilidad del pipeline de generación y entrenamiento es importante.

\section{La interacción búsqueda--red}

La relación fundamental puede escribirse como:

\[
\text{red}
\rightarrow
\text{búsqueda}
\rightarrow
\text{objetivo}
\rightarrow
\text{red}.
\]

Más explícitamente:

\[
f_{\theta_k}
\rightarrow
\mathrm{MCTS}
\rightarrow
\{(s,\pi,z)\}
\rightarrow
\theta_{k+1}.
\]

La búsqueda actúa como un operador de mejora:

\[
\mathcal{I}(f_\theta).
\]

El entrenamiento actúa como un operador de compresión:

\[
\mathcal{C}(\pi,z)
\rightarrow
f_{\theta'}.
\]

Puede verse a la red como una forma de comprimir en parámetros generales el razonamiento realizado por muchas búsquedas anteriores.

\section{Por qué el valor terminal puede afectar decisiones lejanas}

Supóngase una trayectoria:

\[
s_0
\rightarrow
s_1
\rightarrow
s_2
\rightarrow
\cdots
\rightarrow
s_T.
\]

Con resultado final:

\[
z=-1.
\]

El mismo resultado se utiliza como objetivo para múltiples posiciones:

\[
(s_0,\pi_0,-1),
\]

\[
(s_1,\pi_1,+1),
\]

\[
(s_2,\pi_2,-1),
\]

etc., según la perspectiva del jugador al turno.

Esto permite que una consecuencia negativa lejana modifique el valor aprendido de posiciones anteriores.

Sin embargo, esa propagación no es exacta ni instantánea. Depende de:

\begin{itemize}
    \item frecuencia con que reaparezcan estados similares;
    \item capacidad de la red;
    \item distribución de entrenamiento;
    \item calidad de MCTS;
    \item tasa de aprendizaje;
    \item interferencia con otros ejemplos.
\end{itemize}

\section{Diferencia entre retropropagación MCTS y backpropagation neuronal}

Conviene no confundir dos procesos.

\subsection{Backup en MCTS}

Propaga un valor a través del árbol:

\[
v
\rightarrow
Q(s,a).
\]

\subsection{Backpropagation neuronal}

Propaga derivadas:

\[
\nabla_\theta \mathcal{L}
\]

a través de las capas de la red.

Son mecanismos completamente diferentes.

\section{No existe una garantía práctica de optimalidad}

En teoría, ciertos algoritmos MCTS tienen propiedades asintóticas bajo supuestos específicos. Pero AlphaZero real utiliza:

\begin{itemize}
    \item un número finito de simulaciones;
    \item una red aproximadora;
    \item búsqueda selectiva;
    \item entrenamiento estocástico;
    \item datos generados por políticas cambiantes.
\end{itemize}

Por tanto, en una posición concreta no puede afirmarse:

\[
a_{\mathrm{AlphaZero}}
=
a^*.
\]

Lo correcto es afirmar que AlphaZero intenta aproximar una política fuerte mediante mejora iterativa.

\section{Error de aproximación}

Podemos descomponer conceptualmente el error total en varias fuentes:

\[
\epsilon_{\mathrm{total}}
\approx
\epsilon_{\mathrm{red}}
+
\epsilon_{\mathrm{search}}
+
\epsilon_{\mathrm{data}}
+
\epsilon_{\mathrm{optimization}}
+
\epsilon_{\mathrm{sampling}}.
\]

Donde:

\begin{itemize}
    \item \(\epsilon_{\mathrm{red}}\): limitación de representación;
    \item \(\epsilon_{\mathrm{search}}\): búsqueda insuficiente;
    \item \(\epsilon_{\mathrm{data}}\): cobertura incompleta;
    \item \(\epsilon_{\mathrm{optimization}}\): entrenamiento imperfecto;
    \item \(\epsilon_{\mathrm{sampling}}\): ruido estadístico.
\end{itemize}

\section{Presupuesto de búsqueda}

Sea:

\[
M
\]

el número de simulaciones por jugada.

En general, aumentar:

\[
M
\]

mejora la calidad de la búsqueda, pero incrementa el costo:

\[
C_{\mathrm{move}}
\propto
M
\cdot
C_{f_\theta}.
\]

No obstante, la mejora no es lineal. Normalmente existe rendimiento decreciente.

\section{Uso de batching}

Las evaluaciones neuronales pueden agruparse:

\[
\{s_1,\ldots,s_B\}
\rightarrow
f_\theta.
\]

Esto aprovecha hardware paralelo.

El problema es que MCTS es secuencial en parte:

\[
\text{selección}
\rightarrow
\text{evaluación}
\rightarrow
\text{backup}.
\]

La implementación eficiente requiere coordinar muchas búsquedas simultáneas.

\section{Virtual loss}

Cuando varias simulaciones se ejecutan en paralelo, pueden seleccionar la misma rama. Una técnica frecuente es introducir temporalmente una penalización:

\[
Q_{\mathrm{temp}}(s,a)
=
Q(s,a)-L_{\mathrm{virtual}}.
\]

Esto reduce colisiones entre hilos.

\section{Entrenamiento síncrono y asíncrono}

Un sistema distribuido puede contener:

\[
\text{self-play workers},
\]

\[
\text{learner},
\]

\[
\text{parameter server}.
\]

Los workers generan partidas con una versión:

\[
\theta_k.
\]

El learner produce:

\[
\theta_{k+1}.
\]

Un riesgo sistémico es el \emph{policy lag}: partidas nuevas pueden estar siendo generadas todavía con una política vieja.

Si la tasa de generación supera ampliamente la tasa de entrenamiento:

\[
R_{\mathrm{selfplay}}
\gg
R_{\mathrm{learn}},
\]

el buffer acumula experiencias obsoletas.

Este fenómeno es especialmente relevante en implementaciones prácticas.

\section{Distribución de visitas como mejora de política}

Puede interpretarse que MCTS implementa un operador:

\[
\mathcal{M}_M:
(p_\theta,v_\theta)
\mapsto
\pi_M.
\]

Cuando \(M\) aumenta, la política de visitas incorpora más información de búsqueda.

El entrenamiento intenta resolver:

\[
p_{\theta'}
\approx
\pi_M.
\]

Así, la red intenta absorber parte del conocimiento computado por la búsqueda.

\section{Interpretación como búsqueda amortizada}

Supóngase que una búsqueda costosa encuentra repetidamente patrones de buenas decisiones.

El entrenamiento intenta codificarlos en:

\[
\theta.
\]

Esto puede interpretarse como \emph{amortized search}:

\[
\text{búsquedas costosas pasadas}
\rightarrow
\text{predicciones baratas futuras}.
\]

La red funciona como memoria generalizada de muchos procesos de búsqueda.

\section{Relación entre política y valor}

Las dos cabezas están relacionadas pero no son redundantes.

La política responde:

\[
\text{¿dónde buscar?}
\]

El valor responde:

\[
\text{¿qué tan buena es la posición?}
\]

PUCT utiliza:

\[
P
\]

para priorizar expansión y

\[
Q
\]

para incorporar resultados acumulados.

Una política excelente con un valor pobre puede orientar bien inicialmente pero evaluar mal las hojas.

Un valor excelente con una política pobre puede evaluar correctamente posiciones pero desperdiciar búsquedas en ramas irrelevantes.

\section{Aperturas}

AlphaZero no necesita un libro de aperturas explícito.

Si determinadas secuencias producen buenos resultados, las probabilidades:

\[
p_\theta(a\mid s)
\]

y las visitas:

\[
N(s,a)
\]

tienden a concentrarse en ellas.

El ``libro'' queda implícitamente distribuido entre:

\[
\theta
\]

y el comportamiento de MCTS.

\section{Finales}

En finales, la calidad depende de:

\begin{itemize}
    \item cobertura del autojuego;
    \item capacidad de la red;
    \item profundidad efectiva de MCTS.
\end{itemize}

AlphaZero no posee necesariamente una tabla de finales exacta.

Por tanto, puede aproximar muy bien finales frecuentes sin resolver matemáticamente todos ellos.

\section{Estados transpuestos}

En juegos como ajedrez, diferentes secuencias pueden llegar a la misma posición:

\[
s_i=s_j.
\]

Una implementación puede aprovechar:

\[
\text{transposition tables}
\]

o caches de evaluación.

El AlphaZero conceptual se describe normalmente como árbol, aunque implementaciones prácticas pueden compartir información entre estados idénticos.

\section{Simetrías}

En juegos con simetrías, un estado puede transformarse:

\[
s
\rightarrow
g(s)
\]

sin alterar la estructura estratégica esencial.

Puede utilizarse augmentación:

\[
(s,\pi,z)
\rightarrow
(g(s),g(\pi),z).
\]

En ajedrez las simetrías útiles son más restringidas que en Go debido a características como color, orientación de peones y reglas especiales.

\section{Calibración del valor}

Idealmente:

\[
v_\theta(s)
\approx
\E[z\mid s].
\]

Pero una red puede estar bien ordenada sin estar perfectamente calibrada.

Por ejemplo:

\[
v(s_1)>v(s_2)
\]

puede ser correcto aunque sus magnitudes no representen probabilidades exactas.

La fuerza de juego depende tanto del orden relativo como de la calibración.

\section{Política multimodal}

Si varias jugadas son equivalentes:

\[
Q(s,a_1)
\approx
Q(s,a_2),
\]

MCTS puede distribuir visitas entre ellas.

La política objetivo puede ser:

\[
\pi(a_1\mid s)=0.45,
\qquad
\pi(a_2\mid s)=0.50,
\]

en lugar de imponer una única jugada.

Esto proporciona una señal de aprendizaje más rica que una etiqueta discreta.

\section{Crédito temporal}

Una dificultad fundamental del aprendizaje por refuerzo es asignar crédito a decisiones antiguas.

AlphaZero lo aborda indirectamente mediante:

\begin{enumerate}
    \item resultado terminal \(z\);
    \item entrenamiento del valor en cada estado;
    \item búsqueda que explora consecuencias;
    \item repetición de patrones a lo largo de muchas partidas.
\end{enumerate}

No existe una garantía de que una mala decisión aislada sea identificada perfectamente después de una sola partida.

\section{Una formulación compacta del algoritmo}

\subsection*{Inicialización}

Inicializar:

\[
\theta_0.
\]

\subsection*{Iteración \(k\)}

Para múltiples partidas:

\[
s_0\leftarrow\text{posición inicial}.
\]

Mientras la partida no termine:

\[
(\bm{p},v)=f_{\theta_k}(s_t),
\]

\[
\pi_t=\mathrm{MCTS}(s_t,f_{\theta_k}),
\]

\[
a_t\sim \pi_t,
\]

\[
s_{t+1}=T(s_t,a_t).
\]

Cuando termina:

\[
z=\operatorname{resultado}.
\]

Guardar:

\[
(s_t,\pi_t,z_t).
\]

Después actualizar:

\[
\theta_{k+1}
=
\argmin_\theta
\E_{(s,\pi,z)\sim\mathcal{B}}
\left[
(z-v_\theta(s))^2
-
\pi^\top\log p_\theta(s)
+
\lambda\|\theta\|^2
\right].
\]

\section{Pseudocódigo conceptual}

\begin{verbatim}
initialize neural network f_theta
initialize replay buffer B

repeat:
    for each self-play game:
        s = initial_state
        trajectory = []

        while s is not terminal:
            pi = MCTS(s, f_theta)
            trajectory.append((s, pi))
            a ~ pi
            s = next_state(s, a)

        z = game_result(s)

        for (s_t, pi_t) in trajectory:
            B.add(s_t, pi_t, result_from_player_perspective(z, s_t))

    sample minibatches from B

    update theta minimizing:
        (z - v_theta(s))^2
        - pi^T log p_theta(s)
        + lambda ||theta||^2
\end{verbatim}

\section{Qué significa que AlphaZero ``aprende estrategia''}

No existe una variable interna llamada ``estrategia''.

El conocimiento estratégico queda distribuido en los parámetros:

\[
\theta.
\]

Puede manifestarse en:

\begin{itemize}
    \item valoración de estructura;
    \item preferencia por actividad;
    \item seguridad del rey;
    \item valoración de espacio;
    \item sacrificios posicionales;
    \item coordinación de piezas;
    \item transición favorable a finales.
\end{itemize}

Estos conceptos pueden emerger sin estar etiquetados explícitamente.

\section{Conocimiento implícito}

Una representación interna puede contener una característica latente:

\[
h_j(s)
\]

que correlacione con un concepto humano como seguridad del rey.

Pero la red no necesita nombrar ese concepto.

El conocimiento es distribuido:

\[
\theta
=
\{\theta_1,\theta_2,\ldots,\theta_n\}.
\]

Esto dificulta la auditabilidad.

\section{Auditabilidad}

Una debilidad práctica de AlphaZero es que la explicación causal de una decisión no es inmediata.

Puede observarse:

\[
P(s,a),
\quad
Q(s,a),
\quad
N(s,a),
\]

pero no siempre puede reconstruirse un razonamiento humano equivalente.

Para auditar aprendizaje conviene registrar:

\begin{itemize}
    \item política previa;
    \item visitas por acción;
    \item valores \(Q\);
    \item valor de red;
    \item versión del modelo;
    \item profundidad efectiva;
    \item número de nodos;
    \item resultado posterior.
\end{itemize}

\section{Diagnóstico de aprendizaje}

Un sistema puede parecer estar aprendiendo porque la pérdida disminuye:

\[
\mathcal{L}\downarrow
\]

sin mejorar realmente su fuerza.

Por ello deben medirse por separado:

\begin{enumerate}
    \item precisión de política;
    \item error de valor;
    \item calibración;
    \item fuerza contra referencias estables;
    \item estabilidad entre versiones;
    \item diversidad del autojuego.
\end{enumerate}

\section{Sobreajuste al autojuego}

Si el sistema juega siempre contra sí mismo, puede explotar debilidades particulares de su propia distribución.

Formalmente aprende sobre:

\[
d_{\pi_\theta}.
\]

Pero una política externa induce:

\[
d_{\pi_{\mathrm{opp}}}.
\]

Si ambas distribuciones difieren mucho, puede aparecer un problema de generalización.

\section{Evaluación externa}

Es conveniente utilizar oponentes de referencia:

\[
\mathcal{O}_1,\mathcal{O}_2,\ldots
\]

y posiciones fijas:

\[
\mathcal{S}_{\mathrm{test}}.
\]

Así puede medirse si una nueva versión mejora de forma general o sólo se adapta a su propio estilo.

\section{Catastrophic forgetting}

El entrenamiento continuo puede mejorar unas regiones del espacio de estados y empeorar otras.

Esto se denomina \emph{catastrophic forgetting}.

Un replay buffer suficientemente diverso ayuda a mitigarlo, pero introduce el problema opuesto de obsolescencia.

Por tanto existe una tensión:

\[
\text{recencia}
\quad\text{vs}\quad
\text{diversidad histórica}.
\]

\section{Parámetros esenciales}

Entre los parámetros más importantes se encuentran:

\begin{longtable}{>{\raggedright\arraybackslash}p{0.28\textwidth}
                  >{\raggedright\arraybackslash}p{0.62\textwidth}}
\toprule
\textbf{Parámetro} & \textbf{Función} \\
\midrule
\(M\) & Número de simulaciones MCTS por jugada.\\
\(c_{\mathrm{puct}}\) & Balance entre exploración y explotación.\\
\(\tau\) & Temperatura de la política de visitas.\\
\(\varepsilon\) & Peso del ruido Dirichlet.\\
\(\alpha\) & Concentración del ruido Dirichlet.\\
\(\eta\) & Tasa de aprendizaje.\\
\(\lambda\) & Regularización \(L_2\).\\
\(B\) & Tamaño de mini-batch.\\
\(|\mathcal{B}|\) & Tamaño del replay buffer.\\
\bottomrule
\end{longtable}

\section{Efecto del número de simulaciones}

Si el prior de red es incorrecto pero existe suficiente búsqueda, MCTS puede corregirlo.

Con pocas simulaciones:

\[
\pi
\approx
p_\theta.
\]

Con más simulaciones:

\[
\pi
\]

puede desviarse sustancialmente de:

\[
p_\theta.
\]

Por eso el número de simulaciones controla parcialmente cuánto ``razonamiento adicional'' se realiza sobre la red.

\section{Profundidad versus anchura}

AlphaZero no utiliza una profundidad fija tradicional.

MCTS distribuye recursos de forma adaptativa.

Una variante prometedora puede recibir:

\[
N_1\gg N_2.
\]

Así, la profundidad efectiva depende del interés de la rama.

Esto es conceptualmente diferente de:

\[
\text{minimax depth-}d.
\]

\section{Comparación con minimax}

Minimax clásico busca:

\[
V_d(s)
=
\max_a
\min_b
\cdots
\]

hasta profundidad fija \(d\), utilizando una función heurística en las hojas.

AlphaZero reemplaza:

\begin{itemize}
    \item profundidad fija por expansión selectiva;
    \item heurística manual por red neuronal;
    \item ordenamiento manual por política aprendida.
\end{itemize}

\section{Comparación con Q-learning}

Q-learning intenta aproximar:

\[
Q^*(s,a).
\]

AlphaZero modela principalmente:

\[
p(a\mid s),
\qquad
v(s).
\]

MCTS genera estimaciones:

\[
Q_{\mathrm{MCTS}}(s,a).
\]

Una ventaja es separar:

\[
\text{prior de acción}
\]

de:

\[
\text{evaluación del estado}.
\]

\section{Comparación con policy gradient puro}

Un policy gradient puro actualiza directamente:

\[
\nabla_\theta
\E[z].
\]

AlphaZero utiliza una política objetivo más informativa:

\[
\pi_{\mathrm{MCTS}}.
\]

Por ello cada estado contiene una señal estructurada sobre múltiples acciones, no solamente el resultado terminal.

\section{El papel de MCTS como ``teacher''}

Puede interpretarse MCTS como un maestro generado dinámicamente.

La red actual:

\[
f_{\theta_k}
\]

crea una búsqueda que produce:

\[
\pi_k.
\]

Luego esa política se convierte en objetivo para:

\[
f_{\theta_{k+1}}.
\]

Es una forma de auto-destilación guiada por búsqueda.

\section{Riesgo de auto-confirmación}

Existe un riesgo conceptual:

\[
\text{red sesgada}
\rightarrow
\text{búsqueda sesgada}
\rightarrow
\text{datos sesgados}.
\]

La exploración, ruido, autojuego diverso y suficiente presupuesto de búsqueda son mecanismos que ayudan a romper ese circuito.

\section{Convergencia: qué puede y qué no puede afirmarse}

En un sistema ideal tabular, ciertos esquemas de policy iteration poseen garantías fuertes.

AlphaZero viola muchas hipótesis simples:

\begin{itemize}
    \item aproximación funcional;
    \item entrenamiento estocástico;
    \item búsqueda limitada;
    \item distribución cambiante;
    \item optimización no convexa.
\end{itemize}

Por tanto, no existe una garantía general práctica de:

\[
\theta_k
\rightarrow
\theta^*
\]

ni de:

\[
\pi_{\theta_k}
\rightarrow
\pi^*.
\]

El éxito se establece empíricamente.

\section{Una lectura sistémica}

AlphaZero puede verse como un sistema de tres niveles:

\[
\boxed{
\text{red}
\rightarrow
\text{búsqueda}
\rightarrow
\text{aprendizaje}
}
\]

La red aporta velocidad.

La búsqueda aporta deliberación.

El aprendizaje convierte deliberación pasada en conocimiento reutilizable.

La calidad total depende del equilibrio entre los tres.

\section{Implicancias para implementaciones experimentales}

Para un motor experimental que busque inspiración en AlphaZero, conviene preservar conceptualmente las siguientes separaciones:

\begin{enumerate}
    \item \textbf{evaluación de posición};
    \item \textbf{prior de jugadas};
    \item \textbf{búsqueda};
    \item \textbf{generación de experiencia};
    \item \textbf{aprendizaje};
    \item \textbf{persistencia};
    \item \textbf{evaluación independiente}.
\end{enumerate}

Mezclar estos niveles dificulta diagnosticar si una mejora proviene de:

\[
\text{mejor red},
\quad
\text{más búsqueda},
\quad
\text{mejores datos},
\]

o simplemente:

\[
\text{más costo computacional}.
\]

\section{Qué debería medirse}

Una implementación auditable debería registrar al menos:

\[
\text{policy entropy},
\]

\[
\text{value error},
\]

\[
\text{search improvement},
\]

\[
\KL(\pi_{\mathrm{MCTS}}\parallel p_\theta),
\]

\[
\text{win rate},
\]

\[
\text{draw rate},
\]

\[
\text{search nodes},
\]

\[
\text{latency},
\]

\[
\text{model version}.
\]

Una métrica particularmente útil es:

\[
\Delta_{\mathrm{search}}
=
V_{\mathrm{after\ search}}
-
V_{\mathrm{network}}.
\]

Aunque su definición concreta requiere cuidado, conceptualmente mide cuánto valor adicional aporta la búsqueda.

\section{Desacople entre generación y aprendizaje}

Si se utilizan muchos workers, la arquitectura puede representarse como:

\[
W_1,W_2,\ldots,W_n
\rightarrow
Q
\rightarrow
L.
\]

Donde:

\begin{itemize}
    \item \(W_i\): workers de autojuego;
    \item \(Q\): cola o replay buffer;
    \item \(L\): learner.
\end{itemize}

La condición deseable no es simplemente maximizar throughput.

Debe controlarse la edad de experiencia:

\[
A
=
k_{\mathrm{current}}
-
k_{\mathrm{generator}}.
\]

Un valor alto de \(A\) indica que se está entrenando o generando con políticas antiguas.

\section{Edad de política y coherencia}

Si una partida comenzó con:

\[
\theta_k,
\]

es conceptualmente limpio que termine con la misma versión o con una política coherente.

Cambiar la red a mitad de una partida produce una trayectoria generada por:

\[
\theta_k,\theta_{k+1},\ldots
\]

lo que dificulta interpretar los datos.

No es necesariamente inválido, pero reduce la trazabilidad experimental.

\section{Conclusiones}

AlphaZero no es simplemente una red neuronal que elige jugadas, ni un MCTS clásico, ni un algoritmo Bellman tabular. Su fuerza proviene de un circuito cerrado:

\[
\text{predicción}
\rightarrow
\text{búsqueda}
\rightarrow
\text{experiencia}
\rightarrow
\text{aprendizaje}.
\]

La red proporciona:

\[
(p,v).
\]

PUCT utiliza estas estimaciones para construir una búsqueda selectiva.

La búsqueda produce:

\[
\pi.
\]

El resultado terminal produce:

\[
z.
\]

El entrenamiento ajusta:

\[
p_\theta\rightarrow\pi,
\qquad
v_\theta\rightarrow z.
\]

El proceso se repite.

La ``mejor estrategia'' no aparece como una fórmula cerrada ni como una tabla exacta. Emerge progresivamente de la interacción entre autojuego, búsqueda y aproximación neuronal.

La conclusión técnicamente más importante es que AlphaZero no elimina la búsqueda: aprende a hacerla más eficiente. Tampoco elimina la necesidad de evaluación: aprende una función de valor. Y tampoco demuestra optimalidad: construye una aproximación cada vez más fuerte dentro de un presupuesto computacional finito.

\section{Resumen matemático final}

La red:

\[
f_\theta(s)
=
(p_\theta(\cdot\mid s),v_\theta(s)).
\]

Selección PUCT:

\[
a
=
\argmax_a
\left[
Q(s,a)
+
c_{\mathrm{puct}}P(s,a)
\frac{\sqrt{\sum_b N(s,b)}}
{1+N(s,a)}
\right].
\]

Política MCTS:

\[
\pi(a\mid s)
=
\frac{N(s,a)^{1/\tau}}
{\sum_b N(s,b)^{1/\tau}}.
\]

Ejemplo de entrenamiento:

\[
(s,\pi,z).
\]

Pérdida:

\[
\mathcal{L}(\theta)
=
(z-v_\theta(s))^2
-
\sum_a
\pi(a\mid s)\log p_\theta(a\mid s)
+
\lambda\|\theta\|_2^2.
\]

Actualización:

\[
\theta
\leftarrow
\theta
-
\eta\nabla_\theta\mathcal{L}.
\]

Ciclo:

\[
f_{\theta_k}
\rightarrow
\mathrm{MCTS}
\rightarrow
(s,\pi,z)
\rightarrow
f_{\theta_{k+1}}.
\]

\appendix

\section{Glosario}

\begin{description}
    \item[Estado \(s\)] configuración completa del juego.
    \item[Acción \(a\)] jugada legal.
    \item[Política \(p\)] distribución de jugadas producida por la red.
    \item[Política \(\pi\)] distribución mejorada producida por MCTS.
    \item[Valor \(v\)] predicción de resultado.
    \item[Prior \(P\)] probabilidad inicial asignada a una acción en MCTS.
    \item[Visitas \(N\)] número de veces que una arista fue explorada.
    \item[Valor medio \(Q\)] evaluación acumulada de una acción.
    \item[PUCT] regla que combina valor y exploración guiada por prior.
    \item[Replay buffer] memoria de experiencias de entrenamiento.
    \item[Self-play] autojuego.
    \item[Policy lag] retraso entre la política actual y la que generó los datos.
    \item[Bootstrap] mejora progresiva usando las propias estimaciones.
\end{description}

\section{Referencias conceptuales sugeridas}

Para profundizar en el tema conviene consultar los trabajos originales de AlphaGo Zero y AlphaZero, además de literatura sobre Monte Carlo Tree Search, aprendizaje por refuerzo, policy iteration, redes residuales y métodos actor--critic.

\end{document}
